\section{おわりに}
\label{sec:owarini}

本稿では，回転機構のレーザー受光と楽器間無線通信による同期・輪唱システムの設計・実装・評価について報告した．本節では，プロジェクトで実現できた内容と残された課題を整理し，\ref{sec:hajimeni}節で示した目的との対応関係を述べる．

\ref{sec:hajimeni}節では，デジタル配信では代替できないライブ体験の価値が，同じ空間を共有しているという感覚に由来するという知見に基づき，既存の自動演奏技術が持つ「多感覚的フィードバック」と「空間全体への共有」の不足という課題を指摘した．そのうえで，ユーザの能動的操作を受け付ける指揮デバイス，空間全体へ視覚的フィードバックを共有する観覧車型の中継機デバイス，および分散配置された楽器デバイスの3種で構成されるシステムにより，会場全体を巻き込む没入感の高い演奏体験を提供することを目的として設定した．

この目的に対し，本プロジェクトで実現できた内容は次の通りである．第一に，能動的操作の実現である．ボタン押下による演奏の開始・終了，指揮デバイスの振り動作および可変抵抗器によるBPM・音量のリアルタイム変更を実装し，\ref{sec:hyouka}節の評価において，振り動作によるBPM・音量変更は反映成功率100\%，ボタン押下は成功率98.8\%・平均応答時間約\SI{1.0}{s}で動作することを確認した．第二に，空間的フィードバックの実現である．観覧車型デバイスの回転とレーザー受光による発音制御を実装し，物理的に分散した複数の楽器デバイスを，一度同期が確立した後はハース効果の許容限界である\SI{50}{ms}未満の演奏ズレで協調動作させた．これは，複数の発音体が聴覚上1つの輪唱として知覚される精度である．第三に，音色再現の実現である．6種の楽器音を合成し，MOS評価に基づく音色再現度のMOP（中央値および最頻値がともに4.0以上）を全楽器で達成した．さらに，本チームの評価とは独立に実施された参加者共通アンケートにおいて，「演出の華やかさ」および「楽しさ・高揚感」が全参加チームの平均を有意に上回ったことは，能動的操作と空間的フィードバックの組み合わせという本システムの構成が，体験として実際に機能したことを裏付けている．

一方，残された課題も評価を通じて明確になった．第一に，同期確立時間の課題である．同期開始条件を\SI{50}{ms}以下とした場合，同期の確立自体に長時間を要し，応答遅延のMOPを満たせない試行が生じた．この要因は，パケットロス対策としての冗長送信と同期確立時間とのトレードオフ，およびモータ由来の電気的ノイズにあると分析した．第二に，体験としての「同期の心地よさ」の課題である．参加者共通アンケートにおいて「一体感・同期の心地よさ」は唯一全体平均を下回り，項目間の比較でも「演出・楽しさ」に対して有意に低い評価にとどまった．没入感についても，平均値としては全体を上回る傾向がみられたものの，順位に基づく検定では有意差を確認できなかった．すなわち，空間的フィードバックの演出面の効果は明確に確認された一方，序論で目的とした没入感の向上は部分的な達成にとどまり，その隘路が同期確立の高速化・安定化にあることが，技術的評価と体験評価の両面から一貫して示された．第三に，実装・評価上の課題として，ブレッドボード実装に起因する操作反映の不安定さ，レーザー光量の電池残量への依存や光学系の位置調整に代表されるハードウェア依存性，ピアノ・キックドラムの音色再現のばらつき，および試行回数と目標値の水準の不整合といった評価設計の限界が残っている．

以上を総括すると，本プロジェクトは，能動的操作と空間全体で共有される物理的フィードバックを兼ね備えた演奏システムを実装し，その演出効果を定量的な評価によって実証した点で，序論で示した目的を大きく前進させた．同時に，没入感の完全な達成には，利用者がいつ体験しても同期した演奏に接することができる同期確立の即時性が不可欠であることを明らかにした．\ref{sec:kousatsu}節で述べた確認応答に基づく再送方式や時刻同期方式への変更，ノイズ対策，基板実装への移行といった改善を施したうえで，同期確立前後の体験を統制した評価を行うことが，会場全体を巻き込む没入感の高い演奏体験の完成に向けた次の段階である．
